\documentclass [varwidth]{standalone}
%scale[1]
%lmargin[30]
%vmargins[30]
\usepackage {texmaths-main}
\begin {document}
\setCompleteVersion {0} \begin {tikzpicture}[scale=0.8] \coordinate (O)at(0,0); \coordinate (O1)at(6,0); \coordinate (-O1)at(-6,0); \coordinate (O2)at(0,5); \coordinate (-O2)at(0,-5); \draw [-latex](-O1)--(O1); \draw [gray,line width=0.1pt](-6,-5)grid(6,5); \draw [-latex](-O2)--(O2); \draw [white,line width=0.5pt](-6,-5)rectangle(6,5); \foreach \x in{-4,...,4}{ \draw [line width=1pt](\x ,-0.1)--(\x ,0.1); \draw [line width=1pt](-0.1,\x )--(0.1,\x ); } \coordinate (A)at(1,3); \coordinate (B)at(4,-2); \coordinate (C)at(-1.5,1.5); \coordinate (D)at(-1.5,-2); \foreach \X in{A,B,C,D}{ \draw node at(\X ){+}; \draw node[above right]at(\X ){\X }; } \draw node[below right]at(O){0}; \draw node[below right]at(1,0){1}; \end {tikzpicture}\vskip 20pt\noindent Le point de coordonées (0;0) se nomme \toComplete {l'origine du repère}.\par \noindent L'axe horizontal est se nomme l'axe des \toComplete {abscisses}.\par \noindent L'axe vertical est se nomme l'axe des \toComplete {ordonnées}.\par \noindent Le point A a pour coordonnées \toComplete {(+1 ; +3)}.\par \noindent Le point B a pour coordonnées \toComplete {(+4 ; -2)}.\par \noindent Le point C a pour coordonnées \toComplete {(-1,5 ; +1,5)}.\par \noindent Le point D a pour coordonnées \toComplete {(-1,5 ; -2)}.
\end {document}
